\documentclass[../../../include/open-logic-section]{subfiles}

\begin{document}

\olfileid{sth}{z}{infinity-again}	
\olsection{Tanpa Wates}

Kita wis duwe aksioma kang cukup kanggo njamin anane
himpunan kang cacahé tanpa wates (yen paling ora ana siji).
Upamane ana sawijining himpunan, mula $\emptyset$ ana
(miturut \olref[sth][z][sep]{prop:emptyexists}). Saiki
kanggo sembarang himpunan~$x$, himpunan $x \cup \{x\}$ ana
miturut \olref[sth][z][pairs]{prop:pairsconsequences}.
Yen proses iki dibaleni sawetara kaping, kita entuk
himpunan-himpunan iki:
\begin{enumerate}
	\item[0.] $\emptyset$ %& $0$ & tataran $0$\\
	\item[1.] $ \{\emptyset\}$ %& $1$ & tataran $1$\\
	\item[2.] $\{\emptyset, \{\emptyset\}\}$ %& $2$ & tataran $2$\\
	\item[3.] $\{\emptyset, \{\emptyset\}, \{\emptyset, \{\emptyset\}\}\}$ %&$3$ & tataran $3$\\
	\item[4.] $\{\emptyset, \{\emptyset\}, \{\emptyset, \{\emptyset\}\}, \{\emptyset, \{\emptyset\}, \{\emptyset, \{\emptyset\}\}\}\}$% & $4$ & tataran $4$ 
\end{enumerate}
lan kita bisa mriksa yen saben himpunan iki beda.

Kita miwiti nomer saka $0$ amarga sawetara alasan.
Salah sijine yaiku iki: ora angel mriksa yen himpunan
kang diwenehi label ``$n$'' nduweni persis $n$ anggota
lan (kanthi intuitif) kawangun ing tataran kaping $n$.

Nanging iki menehi kita himpunan kang cacahé
\emph{tanpa wates}, tanpa njamin anane \emph{himpunan
tanpa wates}, yaiku siji himpunan kanthi anggota kang
cacahé tanpa wates. Iki wigati: tanpa himpunan tanpa
wates (miturut Dedekind), kita ora bisa nggawe aljabar
Dedekind. Kita mbutuhake aljabar Dedekind supaya bisa
nganggep aljabar iku minangka himpunan wilangan asli.
(Bandhingna \olref[sfr][infinite][dedekindsproof]{sec}.)

Mligine, aksioma kang wis kita aturake nganti saiki
\emph{ora} njamin anane himpunan tanpa wates. Mula
kita kudu nambah aksioma anyar:

\begin{axiom}[Tanpa Wates]
Ana himpunan $I$ kanthi $\emptyset \in I$ lan
$x \cup \{x\} \in I$ saben $x \in I$.
\begin{align*}
	\exists I( & (\exists o \in I)\forall x\ x \notin o \land {}\\
	& (\forall x \in I)(\exists s \in I)\forall z(z \in s \liff (z \in x \lor z = x)))
\end{align*}
\end{axiom}

Saka himpunan $I$ kang diwenehake Aksioma Tanpa Wates,
kita bisa njupuk himpunan bagean kang tanpa wates miturut
Dedekind: tutupan saka himpunan kosong miturut fungsi penerus.
Unsur kang dipilih yaiku $\emptyset$, lan tutupan ing
$I$ nduweni injeksi $s(x) = x\cup \{x\}$.
Saiki \olref[sfr][infinite][dedekind]{thm:DedekindInfiniteAlgebra}
nuduhake carane njupuk Aljabar Dedekind saka himpunan
tanpa wates miturut Dedekind; aljabar iki bakal kita
anggep minangka himpunan wilangan asli. Luwih trep:
\begin{defn}\ollabel{defnomega}
Upamane $I$ iku sembarang himpunan kang diwenehake Aksioma
Tanpa Wates. Upamane $s$ iku fungsi $s(x) = x \cup \{x\}$.
Upamane $\omega = \closureofunder{s}{\emptyset}$. Anggota
$\omega$ kita arani \emph{wilangan asli}, lan $n$ diarani
asil saka $n$ penerapan $s$ marang $\emptyset$.
\end{defn}

Saiki kita bisa ndeleng maneh yen himpunan kang diwenehi
label ``$n$'' sawetara paragraf sadurunge bakal dianggep
\emph{minangka} wilangan $n$.

Kita bakal ngrembug teges saka panetepan iki ing
\olref[sth][z][nat]{sec}. Saiki panetepan iki ngidini
kita mbuktekake asil intuitif:

\begin{prop}\ollabel{naturalnumbersarentinfinite}
Ora ana wilangan asli kang tanpa wates miturut Dedekind.
\end{prop}

\begin{proof}
Bukti iki nganggo induksi, yaiku
\olref[sfr][infinite][induction]{thm:dedinfiniteinduction}.
Cetha yen $0 = \emptyset$ ora tanpa wates miturut Dedekind.
Kanggo langkah induksi, kita mbuktekake kontrapositife:
yen (sacara absurd) $s(n)$ tanpa wates miturut Dedekind,
mula $n$ tanpa wates miturut Dedekind.

Dadi upamane $s(n)$ tanpa wates miturut Dedekind,
yaiku ana sawenehing !!{injection} $f$ kanthi
$\ran{f}\subsetneq \dom{f} = s(n) = n \cup \{n\}$.
Ana rong kasus kang kudu ditimbang.

\emph{Kasus 1: $n \notin \ran{f}$.} Mula
$\ran{f} \subseteq n$, lan $f(n) \in n$.
Upamane $g = \funrestrictionto{f}{n}$; banjur
$\ran{g} = \ran{f} \setminus \{f(n)\} \subsetneq n = \dom{g}$.
Mula $n$ tanpa wates miturut Dedekind.

\emph{Kasus 2: $n \in \ran{f}$.} Pilihen
$m \in \dom{f} \setminus \ran{f}$, banjur temtokna
fungsi $h$ kanthi domain $s(n) = n \cup \{n\}$:
\[
h(x) = 
	\begin{cases}
		f(x) & \text{yen }f(x) \neq n\\
		m & \text{yen }f(x)=n
	\end{cases}
\]
Dadi $h$ lan $f$ padha ing kabeh panggonan kajaba
$h(f^{-1}(n)) = m \neq n = f(f^{-1}(n))$.
Amarga $f$ iku !!a{injection}, $n \notin \ran{h}$;
lan $\ran{h} \subsetneq \dom{h} = s(n)$. Saiki $n$
tanpa wates miturut Dedekind kanthi nggunakake
penalaran Kasus 1.
\end{proof}

Nanging isih ana pitakonan carane kita bisa \emph{ndhukung}
Aksioma Tanpa Wates. Wangsulan cekake yaiku kita kudu
nambah prinsip liyane marang gambaran sadurunge.
Prinsip iku mangkene:
\begin{enumerate}
	\item[] \stagesinf. Ana tataran tanpa wates. Yaiku ana
	tataran kang (a) dudu tataran kapisan, (b) nduweni
	sawetara tataran sadurunge, nanging (c) ora nduweni
	pendahulu langsung.
\end{enumerate}
Aksioma Tanpa Wates langsung ndherek saka prinsip iki.
Kita ngerti yen wilangan asli $n$ kawangun ing tataran $n$.
Mula himpunan $\omega$ kawangun ing tataran tanpa wates
kapisan. Lan $\omega$ dhewe dadi saksi Aksioma Tanpa Wates.

Nanging iki mung mbalekake kita marang pitakonan carane
ndhukung \stagesinf. Kaya \stagessucc, prinsip iki ora
kalebu kanthi gamblang ing gambaran hierarki
kumulatif-iteratif sadurunge. Malah ora ana ing gagasan
hierarki iteratif dhewe, kang nggawe himpunan tataran
demi tataran, kang meksa kita nganggep proses iku
nduweni tataran \emph{tanpa wates}. Katon lumrah yen
tataran-tataran iku mung diurutake kaya wilangan asli.

Nanging saka kene tuwuh masalah kang cetha. Ing
\olref[sfr][infinite][dedekindsproof]{sec}, kita wis
nimbang ``bukti'' Dedekind yen ana himpunan tanpa wates
miturut Dedekind (saka pikiran). Bukti iku bisa wae
ora krasa marem. Yen \stagesinf{} ora ``dijibakake''
dening konsepsi iteratif tumrap himpunan (utawa dening
``hukum-hukum pikiran''), kita isih durung nduweni
alasan intrinsik kanggo negesake anane himpunan
tanpa wates miturut Dedekind.

Mesthi isih akeh kang bisa diandharake. Nanging muga-muga
saiki kita bisa wiwit mikir apa kang \emph{dibutuhake}
kanggo ndhukung sawijining aksioma (utawa prinsip).
Ing salajengipun kita mung bakal nganggep \stagesinf{}
lumaku.

\end{document}
